\section{Discussion}\label{sec:discussion}
The main conclusion is a limit on extending polynomial penalty encoding
from linear to convex quadratic native constraints. Finite exactness and
uniform strict feasibility on feasible integer slices do not suffice:
an improving infeasible slice can lie doubly exponentially close to the
linking equality. The upper bounds identify two ways to recover
polynomial encoding, by fixing continuous dimension or the number of
nonlinear native inequalities under the stated regularity assumptions.
These results refine the quantitative scope of existing exactness
theorems; they do not replace their existence theory.

The calibration result then shows why a safe sufficient coefficient and
a nearly smallest one are different computational targets, even when
the original optimizer is supplied. The perturbation result provides a
probabilistic alternative with explicit sampling and feasibility
qualifications. The encoding upper bounds and the perturbation guarantee
allow large numerical coefficients and do not imply fast solution of the
augmented subproblem. Changing the equality, using succinct exponents or
choosing a fractional-power augmentation changes the representation
question. The contributions concern
ordinary binary coefficients for the prescribed norm and residual.

Supplementary scripts check the exact finite envelopes, determinant
identities, degenerate regularization, calibration dual formulas and thresholds, and selected
grid and tail identities with rational or symbolic arithmetic. These
checks support the algebra and indexing; the proofs establish the general
claims. No solver-performance claim is made.
